% =====================================================================
% CHAPTER 6 — MEASUREMENT ERROR AND POSITION SIZING (L3)
%                                            (~18% of total length)
%
% The second novel chapter: G-study -> D-study -> trading rule.
% The arc: decompose the variance, derive per-item reliability, then
% ask whether knowing the reliability is WORTH anything in P/L.
%
% Structure:
%   6.1 Hypotheses (H3a structure / H3b human-tier link / H3c dominance)
%   6.2 G-study: variance decomposition       <- the headline table
%   6.3 Validation against human agreement tiers
%   6.4 D-study: deriving reliability rho_t
%   6.5 Three trading rules compared
%   6.6 Robustness and the publishable null
%   6.7 Answer to RQ3
%
% PITFALLS:
%  - The G-study table (text x model x prompt x sample components) is
%    the first such number for financial LLM sentiment — present it
%    as a RESULT in its own right, not as preprocessing for 6.5.
%  - The three rules must run on the IDENTICAL signal and test window;
%    any divergence and the comparison is dead.
%  - "Weakly dominates" was the registered claim — if the fixed
%    threshold wins, that bounds the economic value of reliability
%    information, and the chapter still stands.
% =====================================================================
\chapter{Measurement Error and Position Sizing (L3)}
\label{ch:layer3}

\section{Hypotheses}
\label{sec:l3-hypotheses}

\begin{itemize}
  \item \textbf{H3a:} sentiment measurement variance is structured ---
        model and prompt facets are material, not noise.
  \item \textbf{H3b:} item-level variance correlates with human
        agreement tiers.
  \item \textbf{H3c:} reliability-aware rules weakly dominate the
        fixed threshold on risk-adjusted P/L.
\end{itemize}

\section{G-Study: Variance Decomposition}
\label{sec:l3-gstudy}

[Crossed design: item $\times$ model $\times$ prompt-variant $\times$
stochastic sample (\cref{sec:data-roster}); variance components via
mixed-effects estimation; framework per
\textcite{cronbach1972, shavelson1991}.]

\begin{table}[htbp]
  \centering
  \caption{Variance components of LLM sentiment measurement.
  [Placeholder --- this table is the chapter's headline result and,
  to our knowledge, the first G-coefficient for financial LLM
  sentiment.]}
  \label{tab:l3-gstudy}
  \begin{tabular}{@{}l S[table-format=2.1]@{}}
    \toprule
    Facet & {Share of variance (\%)} \\
    \midrule
    Item (text)            & 55.0 \\
    Model                  & 20.0 \\
    Prompt variant         &  8.0 \\
    Item $\times$ model    & 10.0 \\
    Stochastic sample      &  4.0 \\
    Residual               &  3.0 \\
    \bottomrule
  \end{tabular}
\end{table}

\section{Validation Against Human Agreement}
\label{sec:l3-validation}

[H3b: do the items that destabilize machine measurement coincide with
the ones humans dispute (PhraseBank tiers)? Either answer is
informative --- shared source of uncertainty, or machine-specific
fragility.]

\section{D-Study: Deriving Reliability}
\label{sec:l3-dstudy}

[From components to per-item and per-day reliability $\rho_t$; show
the time series of $\rho_t$ against news regimes --- when is the
instrument trustworthy?]

\section{Three Rules Compared}
\label{sec:l3-rules}

On the identical signal and test window:
\begin{enumerate}
  \item \textbf{Fixed threshold} (L1): trade when $|S_t| > L$.
  \item \textbf{Shrinkage:} trade on $\tilde{S}_t = \rho_t S_t$, same
        threshold.
  \item \textbf{Signal-to-noise sizing:} position $\propto
        \tilde{S}_t/\sigma_t$ with a volatility-targeting cap; plus an
        abstention variant (no trade when $\rho_t$ below a floor).
\end{enumerate}

[Results table mirroring \cref{tab:l1-results}'s metrics, with and
without costs; bootstrap intervals on Sharpe differences.]

\section{Robustness}
\label{sec:l3-robustness}

[Sensitivity to the variance-component estimation sample; pilot-vs-full
stability; BH correction shared with Chapter 5's family. If rule 1
wins: report the measured upper bound on the economic value of
reliability information --- the registered fallback claim.]

\section{Answer to RQ3}
\label{sec:l3-answer}

[Direct answer mirroring \cref{sec:research-questions}.]
